\chapter{Conclusions}
\label{chap:final-conclusions}

\section{Study Conclusions}

This dissertation set out to develop, evaluate, and rigorously compare a\acrshort{cnn}-based pipeline for metal artifact reduction in\acrshort{ct}, prioritizing anatomical preservation, robustness to domain shift, and reproducibility. The pre-specified primary hypothesis was that reconstruction produces a real, patient-level improvement in structural similarity within the Artifact Zone, while leaving the Clean Zone comparatively unchanged. This hypothesis was confirmed on the internal test partition for both the Single Model and the Dual Model, with correction concentrated where artifacts were present rather than distributed indiscriminately across the image. The pre-specified secondary hypotheses of global\acrshort{psnr},\acrshort{ssim}, and\acrshort{mae} improvement over the corrupted input were likewise confirmed for both architectures. The direct architectural comparison did not identify a decisively superior design. The Single Model showed a small but consistent advantage in correction magnitude, globally and in the Artifact Zone, while the Dual Model showed a small but consistent Clean-Zone preservation advantage. This is a correction-preservation trade-off, not evidence that auxiliary mask supervision helps or hurts reconstruction outright.

The fourth objective, generalization to real external data from\acrshort{tcia}, can only be answered in part. In the absence of a clean reference for these external cases, the internal technical evidence is limited to characterizing the magnitude of the input-to-output change, which was of a plausible order and did not indicate collapse to the input or gross distortion. The completed blinded, paired evaluation by five physicians on 15 external cases (\autoref{sec:external-evatluation-doctors}) adds a further, complementary line of evidence. Readers favored the model output over the artifact input for anatomical preservation, artifact-intensity reduction, and clinical-interpretation preference in a majority of presentations. This preference was most consistent for reduced artifact intensity and least consistent for the harm-flagging item that probes for distortion away from the artifact. However, this panel fell two readers short of the pre-specified target. The 95\% confidence interval for the primary output-versus-input contrast included the no-difference value, and intra-reader repeat agreement was only moderate. This reader evidence is therefore best described as a descriptive, directionally favorable signal from a small, purposely built panel, not as statistically confirmed proof of clinical benefit or of anatomical safety in tissue away from the artifact.

Considering the study as a whole, the technical evidence for internal reconstruction performance, reproducibility, and artifact-targeted structural correction is strong and consistent across seeds and architectures. The evidence for external generalization and perceived clinical benefit is present and directionally favorable, but preliminary. It rests on a small external case set and a small reader panel, and it does not, by itself, establish that either architecture is clinically safe or superior for use on unseen patients. Confirming that claim would require replication with a larger, ideally multi-center reader panel and a larger, more anatomically diverse external case set, as discussed in \autoref{chap:final-discussion}.

\section{Dissertation Contributions}
\begin{itemize}
    \item A patient-level traceability procedure that recovered exam and patient identity absent from the original\acrshort{ctmar} metadata, enabling patient-level partitioning designed to minimize cross-subset information leakage (187/41/41 patients), audited with\acrshort{ms-ssim} and human visual review.
    \item A controlled comparison between a single-supervision and a dual-supervision (auxiliary artifact-mask) U-Net, isolating the effect of the additional supervision under identical training conditions across three deterministic seeds per architecture.
    \item A patient-level, region-of-interest-aware statistical evaluation protocol that separates Artifact-Zone correction from Clean-Zone preservation, using pre-specified hypotheses, Holm-adjusted tests, and bootstrap confidence intervals, together with a fully reproducible pipeline of versioned configurations, environment snapshots, and checkpoint hashes.
    \item A blinded, paired external evaluation protocol combining quantitative input--output characterization on independent\acrshort{tcia} cases with a completed multi-reader clinical assessment, reported together with an explicit account of its scale, precision, and resulting limitations.
\end{itemize}

